\section{Evaluation Context and EMS Gap}
\label{sec:background}

Reviews of air-quality forecasting report substantial use of RMSE, MAE, correlation, and related point-error summaries when comparing statistical and machine-learning models \cite{mendez2023,tang2024,chadalavada2025}. These summaries are informative, but reporting of forecast variability and event sensitivity is less consistent in the audited literature.

Rolling-origin evaluation is established for time-series forecasting because it preserves temporal ordering and limits information leakage \cite{tashman2000,bergmeir2018}. Yet a leakage-free evaluation can still be decisionally one-dimensional if eligibility is determined only by error. The gap addressed here is therefore narrow: evidence is limited on whether explicitly combining error and dynamic-fidelity dimensions materially changes eligibility across sites, model families, and horizons. This is an audit of decision consequences, not a claim that the underlying metrics or the broader literature are new.
